\scanpage{239}
\mspage{185}
% archive sheet number: 238
Il suffirait, il est vrai, de supposer que $W_S=W_S^f$ sur
l'une quelconque des sept catégories de l'étage autres
que $\underline{\mathrm{Parf}}_W\,S$. Cela impliquerait alors que \deleted{pour} les
six autres, on aura $\overline{W_S}=\overline{W_S^f}$ (mêmes saturés
forts dans $\mathrm{Fl}(\unsure{\mathcal{F}|F})$\unsure{,}). \inserted{Un des deux} \unsure{les} plus faibles \deleted{\unsure{de ces}} axiomes
en question serait :

Si $f:X\longrightarrow Y$ est un foncteur \uline{cartésien} de catégories
\inserted{(et $W$-fibrées$^{*}$)} \uline{fibrées} sur une même $S$, si $f\in W$, alors $f_s\in W$
pour \deleted{tout} $s\in S$,
% slanted
\margin{``$W$-fibrées'' signifiant ici que pour $u\in\mathrm{Fl}\,S$, les foncteurs $u^*$ de chgt de base $X_{s'}\mathrel{\unsure{\text{---}|\longrightarrow}}X_s$, $Y_{s'}\mathrel{\unsure{\text{---}|\longrightarrow}}Y_s$ sont dans $W$}
(et \inserted{l'énoncé} la version duale, pour les foncteurs cocartésiens de
catégories cofibrées). \deleted{\unsure{Je crois}} Sauf erreur, cet axiome
doit déjà impliquer que $\overline W$ satisfait la version
$W(4\mathrm a)$ la plus forte de toutes (a fortiori $W(4\mathrm b)$)
des axiomes $W(4)$. Cet axiome me paraît donc à
présent beaucoup plus raisonnable que $W(4\mathrm b')$, et
ceci d'autant plus qu'il se réfère directement à la
description de $\mathrm{HOT}^{\mathrm{lc}}_W(S)$ en termes de foncteurs
(ou de ``pseudofoncteurs'' si on ne veut supposer les catégories
scindées et les foncteurs compatibles aux scindages)
$S^{\unsure{o|0}}\mathrel{\unsure{\longrightarrow|\text{---}}}\mathrm{Cat}$. D'ailleurs, bien qu'il n'en ait pas
l'air, cet axiome \uline{est} autodual, puisque la
version duale, pour $S$ fixée, est la version directe
pour $S^o$ ! Cela rend cet axiome encore plus
appétissant.

Donc pour formuler $W(4\mathrm a')$, j'ai envie
à présent de laisser tomber $\underline{\mathrm{Parf}}_W$, et de faire pareil.
Si on garde deux versions $W(4\mathrm a')$ (l'ancienne)
et $W(4\mathrm a'')$ (la nouvelle), l'une et l'autre ont
la vertu d'impliquer $W(4\mathrm b)$. Et $W(4\mathrm b')$ (nouvelle
version) doit impliquer $W(4\mathrm a)$ (version ultra-forte)
sinon pour $W$, du moins pour $\overline W$.

Mes idées étaient \unsure{sérieusement|drôlement|sûrement|terriblement} embrouillées
